\section{History and Implementation}
Around the year~2000 I was asked to compile a set of coding guardrails for a team of developers that had been new to Java (and most were new to programming, too). That was the very first version of this document, based on and extending the \ccjpl, originally released by SUN\index{SUN} and therefore often referred to as the \bdq{SUN Java Coding Conventions}\index{SUN Java Coding Conventions|see{Code Conventions for the Java Programming Language}}. That first version reflected the way I wrote \Java code at that time as well as several best practices from various sources.

Later versions of this document were created for other developer teams, and several software projects had been delivered successfully, while working with these coding conventions.

Some of the code from these projects is still alive (in 2026), proving the initial statement that coding conventions do help to create maintainable software.

I wrote a lot of code using the rules, guidelines and recommendations from this document (and its predecessors, of course), including some libraries that I published on GitHub\index{GitHub}. Some of the code from these libraries is mentioned in chapter \tqfullref{sec:CodingGuidelines}, starting on page~\pageref{sec:CodingGuidelines}. I also added some stuff as examples in chapter \tqfullref{sec:Examples} in the \bdq{\nameref{sec:Appendices}} part of the document. If you are interested in the whole package, you should have a look to \autocite{TQUADRAT_ORG}.

One word about history: \textit{this} document reflects the \textit{current} best practices, as for the year 2026 and for \Java25. Some of these had been \bdq{best practices} already when Java~1.1 was realesed, others evolved over time, superseding others. And of course this may happen with the rules and recommendations from this document, too. So when reading this in 2030, take with a grain of salt.

%==============================================================================
% End of file!!
%==============================================================================
